% GENERATED FILE. Edit canonical lessons, then run npm run generate:books.
% canonical-source-hash: fnv1a64:9e0af31402ffc5ef
% canonical-lessons: IT-C293-scasso, IT-C293-scassinare, IT-C293-antifurto, IT-C293-allarme, IT-C293-vandalo

\chapter{Ladri in casa}
\label{ch:it-ladri-in-casa}
\hlchaptermodality{293}

\begin{chapteropening}
\textbf{By the end of this chapter:} \emph{I can say that thieves broke into a house: the break-in, the forced window, the burglar alarm and the vandals.}

Complete the last lesson of chapter 293: I can say that thieves broke into a house: the break-in, the forced window, the burglar alarm and the vandals.
\end{chapteropening}

\section[lo scasso]{lo scasso — the break-in, the forced entry}

\label{lesson:IT-C293-scasso}



\begin{quote}
Before the new one: say the Italian for the safe, then the Italian for the loot.
\end{quote}

\subsection*{You'll want to know: lo scasso}
\textbf{lo scasso} — \textquotedblleft{}the break-in, the forced entry\textquotedblright{}.

\begin{quote}
Hanno trovato segni di scasso sulla porta.
\end{quote}

\begin{quote}
\emph{They found signs of forced entry on the door.}
\end{quote}

A police word: \emph{segni di scasso}, signs of forced entry.

\begin{cousinweb}
\emph{scasso} is built on \emph{scassare} \textquotedblleft{}to break open\textquotedblright{}, from \emph{s-} and \emph{cassa} \textquotedblleft{}box\textquotedblright{}, from Latin \emph{capsa} \textquotedblleft{}a case\textquotedblright{} (Lewis \& Short).
\end{cousinweb}

\begin{grammarlens}[title={What you've built}]
One more word for telling, explaining and describing at B1.
\end{grammarlens}

\subsection*{Your turn}
\begin{itemize}
\raggedright
  \item \emph{Say it:} \emph{lo scasso}
  \item \emph{Say it:} \emph{lo scasso}, once more
  \item \emph{Say it:} a sentence of your own with \emph{lo scasso}
  \item \emph{Recall:} say the Italian for the safe, then the Italian for the loot, then say \emph{lo scasso} again
\end{itemize}

\begin{tcolorbox}[breakable,colback=teal!4,colframe=teal!35!black,title={Before you move on}]
What does lo scasso mean? (\textquotedblleft{}The break-in, the forced entry\textquotedblright{}.) Say it once more.
\end{tcolorbox}
\section[scassinare]{scassinare — to force open, to break into}

\label{lesson:IT-C293-scassinare}



\begin{quote}
Before the new one: say the Italian for the loot, then the Italian for the break-in.
\end{quote}

\subsection*{You'll want to know: scassinare}
\textbf{scassinare} — \textquotedblleft{}to force open, to break into\textquotedblright{}.

\begin{quote}
Qualcuno ha scassinato la finestra.
\end{quote}

\begin{quote}
\emph{Someone forced the window.}
\end{quote}

For locks, doors, safes and cars.

\begin{cousinweb}
\emph{scassinare} is built on \emph{scassare} \textquotedblleft{}to break open\textquotedblright{}, from \emph{cassa} \textquotedblleft{}box\textquotedblright{} and Latin \emph{capsa} \textquotedblleft{}a case\textquotedblright{} (Lewis \& Short).
\end{cousinweb}

\begin{grammarlens}[title={What you've built}]
One more word for telling, explaining and describing at B1.
\end{grammarlens}

\subsection*{Your turn}
\begin{itemize}
\raggedright
  \item \emph{Say it:} \emph{scassinare}
  \item \emph{Say it:} \emph{scassinare}, once more
  \item \emph{Say it:} a sentence of your own with \emph{scassinare}
  \item \emph{Recall:} say the Italian for the loot, then the Italian for the break-in, then say \emph{scassinare} again
\end{itemize}

\begin{tcolorbox}[breakable,colback=teal!4,colframe=teal!35!black,title={Before you move on}]
What does scassinare mean? (\textquotedblleft{}To force open, to break into\textquotedblright{}.) Say it once more.
\end{tcolorbox}
\section[l'antifurto]{l'antifurto — the burglar alarm, the anti-theft device}

\label{lesson:IT-C293-antifurto}



\begin{quote}
Before the new one: say the Italian for the break-in, then the Italian for to force open.
\end{quote}

\subsection*{You'll want to know: l'antifurto}
\textbf{l'antifurto} — \textquotedblleft{}the burglar alarm, the anti-theft device\textquotedblright{}.

\begin{quote}
L'antifurto della macchina ha suonato tutta la notte.
\end{quote}

\begin{quote}
\emph{The car alarm went off all night.}
\end{quote}

Also an adjective: \emph{un sistema antifurto}.

\begin{cousinweb}
\emph{antifurto} joins Greek \emph{anti-} \textquotedblleft{}against\textquotedblright{} with \emph{furto} \textquotedblleft{}theft\textquotedblright{}, which is inherited from Latin \emph{fūrtum} (Lewis \& Short).
\end{cousinweb}

\begin{grammarlens}[title={What you've built}]
One more word for telling, explaining and describing at B1.
\end{grammarlens}

\subsection*{Your turn}
\begin{itemize}
\raggedright
  \item \emph{Say it:} \emph{l'antifurto}
  \item \emph{Say it:} \emph{l'antifurto}, once more
  \item \emph{Say it:} a sentence of your own with \emph{l'antifurto}
  \item \emph{Recall:} say the Italian for the break-in, then the Italian for to force open, then say \emph{l'antifurto} again
\end{itemize}

\begin{tcolorbox}[breakable,colback=teal!4,colframe=teal!35!black,title={Before you move on}]
What does l'antifurto mean? (\textquotedblleft{}The burglar alarm, the anti-theft device\textquotedblright{}.) Say it once more.
\end{tcolorbox}
\section[l'allarme]{l'allarme — the alarm}

\label{lesson:IT-C293-allarme}



\begin{quote}
Before the new one: say the Italian for to force open, then the Italian for the burglar alarm.
\end{quote}

\subsection*{You'll want to know: l'allarme}
\textbf{l'allarme} — \textquotedblleft{}the alarm\textquotedblright{}.

\begin{quote}
Qualcuno ha dato l'allarme.
\end{quote}

\begin{quote}
\emph{Someone raised the alarm.}
\end{quote}

\emph{Dare l'allarme} — to raise the alarm.

\begin{cousinweb}
\emph{allarme} comes from the cry \emph{all'arme!} \textquotedblleft{}to arms!\textquotedblright{}; \emph{arme} goes back to Latin \emph{arma} \textquotedblleft{}weapons\textquotedblright{} (Lewis \& Short).
\end{cousinweb}

\begin{grammarlens}[title={What you've built}]
One more word for telling, explaining and describing at B1.
\end{grammarlens}

\subsection*{Your turn}
\begin{itemize}
\raggedright
  \item \emph{Say it:} \emph{l'allarme}
  \item \emph{Say it:} \emph{l'allarme}, once more
  \item \emph{Say it:} a sentence of your own with \emph{l'allarme}
  \item \emph{Recall:} say the Italian for to force open, then the Italian for the burglar alarm, then say \emph{l'allarme} again
\end{itemize}

\begin{tcolorbox}[breakable,colback=teal!4,colframe=teal!35!black,title={Before you move on}]
What does l'allarme mean? (\textquotedblleft{}The alarm\textquotedblright{}.) Say it once more.
\end{tcolorbox}
\section[il vandalo]{il vandalo — the vandal}

\label{lesson:IT-C293-vandalo}



\begin{quote}
Before the new one: say the Italian for the burglar alarm, then the Italian for the alarm.
\end{quote}

\subsection*{You'll want to know: il vandalo}
\textbf{il vandalo} — \textquotedblleft{}the vandal\textquotedblright{}.

\begin{quote}
I vandali hanno rotto tutte le finestre della scuola.
\end{quote}

\begin{quote}
\emph{Vandals broke all the school windows.}
\end{quote}

Plural \emph{i vandali}.

\begin{cousinweb}
\emph{vandalo} comes from Latin \emph{Vandalus} (EtymDB), the name of a Germanic people who sacked Rome in the year 455.
\end{cousinweb}

\begin{grammarlens}[title={What you've built}]
One more word for telling, explaining and describing at B1.
\end{grammarlens}

\subsection*{Your turn}
\begin{itemize}
\raggedright
  \item \emph{Say it:} \emph{il vandalo}
  \item \emph{Say it:} \emph{il vandalo}, once more
  \item \emph{Say it:} a sentence of your own with \emph{il vandalo}
  \item \emph{Recall:} say the Italian for the burglar alarm, then the Italian for the alarm, then say \emph{il vandalo} again
\end{itemize}

\begin{tcolorbox}[breakable,colback=teal!4,colframe=teal!35!black,title={Before you move on}]
What does il vandalo mean? (\textquotedblleft{}The vandal\textquotedblright{}.) Say it once more.
\end{tcolorbox}
